\begin{abstract}
本讲义面向论文《疫情监测中的病例规模收益边界：估计目标、环境随机性与外部验证》的作者。讲义的目标是让作者能够不借助外部材料独立重现论文中的每一步推导，并说清每个数值的来源。讲义先列出前置知识、全部符号和模型假设，随后对命题 1 至命题 6、定义 1 与推论 1 逐一展开。每个结论都按同一顺序处理。先照录原式并注明原文式号，再写出每个参数的含义。然后逐步补全论文跳过的推导，并指明每一步所用的工具。最后用论文图 1 的参数给出数值例，并对应到论文的仿真与实证结果。

讲义的第二部分讨论从理论到实证的映射，涵盖可行广义最小二乘、经验交叉点、州聚类 Bootstrap、工具变量和稳健性检验，并汇总 199 条符号与数值核验的结果：187 条通过、10 条存在偏差、2 条不通过。代数结论全部成立。需要修改的内容集中在近似的方向、检验水平的措辞、离散功效的锯齿，以及截距和斜率与情景隐含值之间的口径差。讲义按原文行号列出这些问题和修改方案，并给出答辩问答与后续研究方向。附录给出负二项分布的 Poisson--Gamma 混合表示、二项抽稀的封闭性、全方差公式，以及 Fisher 信息与 Cram\'er--Rao 下界的完整证明。

\vspace{0.8em}
\noindent\textbf{关键词：}负二项分支过程；二项抽稀；全方差分解；方差交叉点；监测设计；推导核验
\end{abstract}
